\documentclass[11pt]{article}
\usepackage[margin=0.9in]{geometry}
\usepackage{amsmath,amssymb,amsthm,booktabs,mathrsfs}
\usepackage[T1]{fontenc}
\usepackage{lmodern,microtype}
\usepackage[colorlinks=true,linkcolor=blue,urlcolor=blue]{hyperref}
\newcommand{\norm}[1]{\left\lVert#1\right\rVert}
\newtheorem{theorem}{Theorem}[section]
\newtheorem{lemma}[theorem]{Lemma}
\newtheorem{proposition}[theorem]{Proposition}
\newtheorem{corollary}[theorem]{Corollary}
\theoremstyle{remark}
\newtheorem{remark}[theorem]{Remark}
\title{Toward a four-dimensional Yang--Mills construction\\
\large Notebook 36: soft determinants and bounded coarse corrections}
\author{Research working document}
\date{17 October 2026}
\setlength{\emergencystretch}{5em}
\begin{document}
\maketitle
\begin{abstract}
We factor the actual covariant mean-zero determinant of Notebook 35
into a soft scalar/vector pair and two determinants with one color
fiber per final block. Both coarse determinants have uniformly bounded
actual Maxwell coefficients as the block size grows. Their flat gaps
alone would not prove this: we establish rooted physical source bounds,
with all powers of the fine spacing and all fiber dimensions retained.

A new gauge-covariant Hessian estimate controls the actual effective
classical action through its first two background derivatives, uniformly
in preceding depth. Local formal path gauges isolate curvature; an
ordered Laurent factorization retains the shifted lattice components.
This yields the relative vertices needed for the vector correction.
The auxiliary off-shell family is identified with the actual one only
through the proved on-shell affine jets.

The fine soft pair has uniformly bounded short-time and long-time
affine heat contributions. The latter follows from a local cubic
resolvent formula and Schatten bounds that retain the block projector.
Consequently only the intermediate heat interval can supply the
required logarithm.

The resulting reduction retains the full effective action and original
constraint normalization. It does not yet compute the remaining fine
determinant's logarithmic coefficient or establish continuum
Yang--Mills existence.
\end{abstract}

\section*{Conventions and proof boundary}

The primitive odd ratio \(\ell\) is fixed and \(K=\ell^k\).
Current coordinates use counting pairings; their spacing in final
physical units is \(K^{-1}\). Normalization \(\mathcal N_0\) subtracts
the flat value, and \([t^2]\) means a Taylor coefficient.
All actual identities refer to the homogeneous stationary germ
constructed in the preceding notebooks. Exact finite determinant
identities and finite flat spectral estimates are distinguished from
uniform spatial and curvature estimates on a chosen joint infinite
classical structural system. No uniqueness of that system is assumed.

For a trace-normalized Cartan element \(H\), put
\[
 C_H=\sum_{q\in\operatorname{spec}(\operatorname{ad}_H)}q^2,
 \qquad \alpha_g=\frac{C_H}{48\pi^2}.
\]
The target remains the actual coefficient
\(22\alpha_g\log K+O(1)\), uniformly in preceding depth.
The finite-part estimates below advance that comparison; they do not
substitute a narrower objective for the existence problem.



\section{A soft scalar factorization with a bounded coarse correction}
\label{ssc-section}

Set \(\mu=K^{-1}\). All fine scalar and bond pairings use counting
measure. In this section the current links are
\(U(e)=\exp(i\mu A(e))\), so \(A\) is a physical potential in final
units. Use the actual covariant moving-reference block average \(B\)
of Notebook 35, and write
\begin{equation}
 J=K^2B^T,\quad J^TJ=I,\quad \Pi=JJ^T,\quad P=I-\Pi,
 \quad \Delta=D^TD.
 \label{ssc-data}
\end{equation}
The normalized block column \(J_y\) has \(K^4\) entries of size
\(K^{-2}\), each an adjoint transport. Distinct columns have disjoint
site supports. Its background dependence stays inside its own block.
Define
\begin{equation}
 T=\Delta+\mu^2\Pi,\qquad
 M=\mu^2J^TT^{-1}J,\qquad
 \widehat G=P\Delta P+\mu^2\Pi .
 \label{ssc-soft}
\end{equation}
The letter \(M\) in this section denotes this coarse scalar matrix;
preceding homogeneous depth will be written \(M_{\rm old}\).

\begin{theorem}[Exact scalar factorization]
\label{ssc-factorization}
On the actual local germ, continued through flatness,
\begin{equation}
 \det\widehat G=\det T\,\det M,\qquad
 \mathcal N_0\log\det\widehat G
   =\mathcal N_0\log\det T+\mathcal N_0\log\det M .
 \label{ssc-determinant}
\end{equation}
Both factors are gauge covariant and have positive flat spectra.
Their physical complex domain can be chosen independently of \(K\).
Moreover,
\begin{equation}
 \left|\beta(\mathcal N_0\log\det M)\right|\le C
 \label{ssc-coarse-bound}
\end{equation}
uniformly in \(K\), preceding depth, and the chosen joint structural
systems. The coefficient is the actual affine Maxwell coefficient.
\end{theorem}

\subsection{Exact algebra and uniform flat estimates}

Choose an orthonormal frame \(V\) for \(\ker B\) on real backgrounds.
In the frame \([V,J]\), the upper block of \(T\) is
\(G_B=V^T\Delta V\). Schur complementation gives
\[
 \det G_B=\det T\,\det(J^TT^{-1}J).
\]
Since \(\det\widehat G=(\mu^2)^{\dim S_c}\det G_B\), this proves
\eqref{ssc-determinant}. Positivity at flatness fixes the logarithm
branches; holomorphic uniqueness continues the identity. This is
finite determinant algebra, with no separately singular globally
based scalar factor.

Write a subscript \(0\) for flatness. Internal-block Neumann Poincare
gives
\[
 \norm{P_0\phi}^2\le C_P K^2\norm{D_0\phi}^2.
\]
Only internal block edges are used, on which a block constant has
zero derivative. Consequently
\begin{equation}
 c(\Delta_0+\mu^2I)\le T_0\le\Delta_0+\mu^2I,\qquad
 c=\bigl(1+\max(C_P,1)\bigr)^{-1}.
 \label{ssc-ellipticity}
\end{equation}
There is no derivative of a sharp block indicator in this argument.

The product-tent lift from Notebook 35, composed with \(J_0\), gives
a map \(E:S_c\to S\) satisfying
\[
 J_0^TE=I,\qquad \norm E\le C_0,\qquad
 \norm{D_0E}\le C_1\mu .
\]
Let \(S\) be the minimum energy with prescribed \(J_0^T\phi=b\).
The feasible field \(Eb\) proves \(0\le S\le C_1^2\mu^2I\).
The lower inverse block then gives
\begin{equation}
 M_0=(I+\mu^{-2}S)^{-1},\qquad
 (1+C_1^2)^{-1}I\le M_0\le I .
 \label{ssc-coarse-gap}
\end{equation}
At \(K=1\), the same statement uses the identity lift and the bounded
unit-lattice gradient.

\subsection{Uniform physical domain and weighted inverses}

The established ordinary-reference trajectory bounds give
intermediate angles \(O(\ell^j\norm A_\infty/K)\). Their sum over
\(j<k\) is bounded by \(C\norm A_\infty\). The hierarchical
adjoint transports defining \(J\) therefore satisfy, on a fixed
complex physical ball,
\begin{equation}
 \norm{J(A)-J_0}\le C\norm A_\infty .
 \label{ssc-J-bound}
\end{equation}
This is an operator-valued block bound: normalizing the \(K^4\)
entries by \(K^{-2}\) removes the counting dimension. Block Banach
Cauchy gives the same uniform first-two derivative bounds.
Only the physical supremum norm is needed for this scalar statement.

The local gauge generator obeys
\(\norm{D(A)-D_0}\le C\mu\norm A_\infty\). Expanding its product
and using \eqref{ssc-ellipticity} gives
\begin{equation}
 \norm{T_0^{-1/2}(T(A)-T_0)T_0^{-1/2}}
       \le C(\norm A_\infty+\norm A_\infty^2).
 \label{ssc-relative}
\end{equation}
The penalty difference is \(\mu^2(\Pi(A)-\Pi_0)\), with the same
bound. This is a complex operator-norm estimate relative to positive
real \(T_0\), not an assertion that a complex transpose product is
positive. A Neumann inverse and \eqref{ssc-J-bound} imply
\[
 \norm{T(A)^{-1}}\le C\mu^{-2},\qquad
 \norm{M(A)-M_0}\le C\norm A_\infty .
\]
Together with \eqref{ssc-coarse-gap}, these give a common radius
independent of \(K\).

Conjugate by \(e^{\theta f(x/K)}\), where \(f\) is real and
one-Lipschitz. Use initial-site weights on bonds and center weights
on coarse sites. The covariant difference changes by \(O(\theta\mu)\);
a \(J\) column has bounded physical diameter. The conjugated
transpose carries the opposite weight. The same relative Neumann
argument, with \(C(\theta+\norm A_\infty)\), proves for \(R=T^{-1}\)
\begin{equation}
 \norm R_\theta\le C\mu^{-2},\qquad
 \norm{DR}_\theta+\norm{RD^T}_\theta\le C\mu^{-1},
 \qquad \norm{DRD^T}_\theta\le C .
 \label{ssc-weighted-energy}
\end{equation}
For example, place \(T_0^{-1/2}\) on the two sides of the Neumann
inverse and use \(\norm{D_\theta T_0^{-1/2}}\le C\).
The bounds hold between restricted final blocks, with an extra
\(e^{-c\,d(x,y)}\). Thus
\begin{equation}
 \norm{M_{xy}}\le Ce^{-c\,d(x,y)}.
 \label{ssc-coarse-decay}
\end{equation}
The coarse fiber has fixed color dimension. Schur summation in
\eqref{ssc-coarse-decay} shows that a small coarse exponential
conjugation changes \(M_0\) by \(O(\theta)\). Weighted \(M(A)-M_0\)
is \(O(\norm A_\infty)\). A common logarithm contour and its
resolvents therefore give uniform exponential kernels for \(\log M\).

\subsection{Rooted derivatives and the actual affine coefficient}

Treat physical directions in one final block as one supremum-norm
Banach space. A local direction \(a_j\) satisfies
\[
 \norm{\delta_jD}\le C\mu\norm{a_j},\qquad
 \norm{\delta_jJ}\le C\norm{a_j}.
\]
The second primitive derivatives have the same respective powers
or better and joint local support. Differentiate
\[
 \delta R=-R(\delta T)R,\qquad
 \delta T=(\delta D)^TD+D^T\delta D+
       \mu^2\bigl((\delta J)J^T+J(\delta J)^T\bigr).
\]
Every term in the first-two derivatives of
\(M=\mu^2J^TRJ\) has balanced powers of \(\mu\), by
\eqref{ssc-weighted-energy}. For example,
\[
 \mu^2(J^TRD^T)(\delta D)(RJ)
 \quad\hbox{costs}\quad
 \mu^2\mu^{-1}\mu\mu^{-2}=1.
\]
At two insertions, the intervening pieces are \(R,DR,RD^T\),
or \(DRD^T\), with exactly the powers in
\eqref{ssc-weighted-energy}. Differentiating \(J\) gives local
endpoint terms; a penalty insertion supplies \(\mu^2\).
The number of terms is independent of \(K\).

Insert block projections at the differentiated factors. The
restricted estimates connect every source and matrix endpoint.
For \(n=1,2\), after reducing the exponent,
\begin{equation}
 \norm{\delta_{j_1}\cdots\delta_{j_n}M_{xy}}
 \le C_n\prod_{\ell=1}^n\norm{a_{j_\ell}}\,
     e^{-c\,\tau(x,y,j_1,\ldots,j_n)} .
 \label{ssc-rooted}
\end{equation}
Here \(\tau\) is minimum tree length on the final-block graph.
The differentiated contour formula for \(\log M\) preserves this
estimate: its intermediate sums are over fixed-dimensional coarse
fibers, and an exponential reserve pays them. This proves rooted
second-derivative bounds with every fixed polynomial source moment.
It does not sum individual fine-coordinate derivative norms.

By covariance, \(\rho_y=\mathcal N_0\operatorname{tr}
(\log M)_{yy}\) is individually gauge invariant, and its linear
color jet vanishes. At each fixed \(K\), the actual stationary
polynomial-jet theorem identifies its affine source with
\[
 A_\alpha(x)=t(x_\nu/K)H,\qquad \alpha\ne\nu ,
\]
up to a fine gauge. Its fine link angle is \(t x_\nu H/K^2\).
In source block \(j\) the coefficient of \(t\) has norm at most
\(C(1+|j|)\). Inserting these directions in the second version of
\eqref{ssc-rooted} for \(\log M\) gives a uniformly finite sum.
The invariant-density Ward identity of Notebook 35 identifies
\(\beta=2[t^2]\rho_0\), proving \eqref{ssc-coarse-bound}.
This uses the fixed-\(K\) actual identification and then estimates
the explicit representative; it assumes no uniform stationary
response as \(K\) grows.

It follows that
\begin{equation}
 \beta(\mathcal N_0\log\det\widehat G)
       =\beta(\mathcal N_0\log\det T)+O(1).
 \label{ssc-scalar-reduction}
\end{equation}
Equality of total finite determinants and localized Hessians is
used here. Their chosen diagonal densities need not agree pointwise.
The remaining magnetic comparison for \(T\) is separate: a
dimension-blind Cauchy bound on its fine trace can still grow with
\(K\), and does not compute its logarithmic coefficient.


\section{Covariance and physical relative Hessian vertices}
\label{chr-section}

This section controls the first two background coefficients of an
actual action Hessian.  It does not estimate the moving hard
mean-zero projection, nor compute a logarithmic determinant
coefficient.  The argument uses exact gauge covariance before
taking absolute values.

\subsection{Hypotheses and conclusion}
\label{chr-hypotheses}

Work on the current unit lattice with counting inner products and a
fixed finite bond/color fiber.  Let \(L(a)\) be a transpose-symmetric
family of left-relative bond Hessians.  Assume that it is exactly
covariant under sitewise adjoint frames and that its current-coordinate
Taylor kernels through background degree two at zero have uniform
exponentially rooted absolute row bounds.  The transposed bounds are
included.  For an action Hessian, these are the action kernels through
degree four together with the local relative-coordinate changes.

Write \(q=L(0)\).  Assume that \(q\) is translation invariant, scalar
on each color factor, and exponentially summable, uniformly in the
old classical depth.  In source-site coordinates impose the
componentwise moment conditions
\begin{equation}
 \sum_r q_{\alpha\beta}(r)=0,\qquad
 \sum_r r_\rho q_{\alpha\beta}(r)=0
 \quad\text{for every }\alpha,\beta,\rho .
 \label{chr-moments}
\end{equation}
Thus its Fourier symbol vanishes to order two.  No symmetry of
\(q_{\alpha\beta}(r)\) in \(\alpha,\beta\) at a fixed \(r\) is assumed.
Only the full transpose relation
\(q_{\alpha\beta}(r)=q_{\beta\alpha}(-r)^T\) is used.

The actual application here is on a chosen infinite structural system,
where these continuous-momentum moments have been established.
Explicit periodizations of that system are also covered with the
specified lifted paths.  An arbitrary actual finite-regulator winding
kernel is not silently identified with such a periodization.

Set \(\mu=K^{-1}\leq1\).  For a complete physical field \(A\), define
\begin{equation}
 \|A\|_{\mathrm{phys},1}
 =\|A\|_\infty+
   \mu^{-1}\max_{\substack{e,e'\text{ nearest}\\
                          \text{same component}}}
                     |A(e')-A(e)| .
 \label{chr-physical-norm}
\end{equation}
Let \(L_j=[t^j]L(t\mu A)\), \(j=1,2\), and suppose
\(\|A\|_{\mathrm{phys},1}\leq a_*\).
Here \(D_0\) denotes the ordinary componentwise gradient of a
\emph{bond field}, not the scalar gauge generator.  Put
\(R_0=D_0^TD_0+\mu^2I\).

\begin{theorem}
\label{chr-main}
Under these hypotheses there are operators \(Z_j,E_j\) such that
\begin{equation}
 L_j=D_0^TZ_j+Z_j^TD_0+E_j,\qquad j=1,2,
 \label{chr-decomposition}
\end{equation}
with
\begin{align}
 \|Z_1\|&\leq C\mu a_*,
 &\|E_1\|&\leq C\mu^2a_*,
 \notag\\
 \|Z_2\|&\leq C\mu^2a_*^2,
 &\|E_2\|&\leq C\mu^2a_*^2 .
 \label{chr-bounds}
\end{align}
Consequently
\begin{equation}
 \|R_0^{-1/2}L_jR_0^{-1/2}\|
       \leq C_j a_*^j,\qquad j=1,2.
 \label{chr-relative}
\end{equation}
The constants are independent of \(K\) and of the old depth.
Mixed second coefficients follow by polarization.  For directions
supported in fixed smooth physical collars, the factors admit
exponentially rooted final-block operator bounds with the same
powers of \(\mu\).
\end{theorem}

For the actual homogeneous classical action \(S_M\), the hypotheses
come from the established uniform support and Cauchy bounds, gauge
covariance, and the flat normalization
\[
 q_{M,\mathrm{mid}}(p)=|p|^2I-pp^T+O(|p|^4).
\]
The analytic midpoint-to-source-site phases preserve the zero and
first jets, so they imply \eqref{chr-moments}.  This last input is
stronger than row-sum cancellation alone.

\subsection{Formal local path gauges}
\label{chr-local-gauge}

For each displacement \(r\), choose a translation-covariant
coordinate path of length \(|r|_1\), with the negative-displacement
path defined by reversal.  Denote its adjoint transport by
\(W_{xy}(a)\).  It obeys \(W_{yx}=W_{xy}^T\) and is complex
orthogonal coefficientwise.

Fix the source site \(x\) of a matrix row.  The transports from \(x\)
to every \(z\) define a formal site gauge \(g_x(z)\), with
\(g_x(x)=I\).  A transformed edge is a loop consisting of the two
selected paths and that edge.  Its length is at most
\(C(1+|z-x|)\); a plaquette filling uses at most
\(C(1+|z-x|)^2\) plaquettes.  Write its link-angle jet as
\[
 b(t,z)=t b_1(z)+t^2b_2(z)+O(t^3).
\]
Discrete Stokes at first order and the physical first-difference
bound give
\begin{equation}
 |b_1(z)|\leq C\mu^2a_*(1+|z-x|)^2.
 \label{chr-b1}
\end{equation}
The second coefficient of the product and its logarithm is a sum
of quadratic pairs of first link coefficients.  There are at most
the square of the loop length such pairs.  Hence
\begin{equation}
 |b_2(z)|\leq C\mu^2a_*^2(1+|z-x|)^2.
 \label{chr-b2}
\end{equation}

These are formal coefficient identities.  The potentially
unbounded gauge \(g_x\) is not placed in a global analytic ball.
One first truncates its polynomial jets far from \(x\), applies
the finite covariance identity, and then passes coefficientwise.
The exponentially rooted bounds absorb the polynomial tails.
In particular this step does not assert a finite-period identity
for an unbounded affine field.

Gauge covariance and \(g_x(x)=I\) give, with a consistent choice
of transport orientation,
\begin{equation}
 L(t\mu A)_{x\alpha,y\beta}
        =L(b(t))_{x\alpha,y\beta}W_{xy}(t\mu A).
 \label{chr-gauge-identity}
\end{equation}
The inverse convention gives the same argument with both
orientations reversed.  Define the dressed flat kernel
\[
 L^P(t)_{x\alpha,y\beta}
     =q_{\alpha\beta}(y-x)W_{xy}(t\mu A).
\]
Using \eqref{chr-b1}--\eqref{chr-b2} and the rooted Taylor kernels,
\begin{align}
 \|[t](L-L^P)\|_{\mathrm{Schur}}
     &\leq C\mu^2a_*,
 \notag\\
 \|[t^2](L-L^P)\|_{\mathrm{Schur}}
     &\leq C\mu^2a_*^2 .
 \label{chr-peierls-error}
\end{align}
At second order the terms are \(L'[b_2]\),
\(\frac12L''[b_1,b_1]\), and the product of \(L'[b_1]\)
with the first transport coefficient.  The latter two have,
respectively, the stronger powers \(\mu^4\) and \(\mu^3\).
All sums are fine current-coordinate Schur sums; no final-block
fiber dimension is inserted.

\subsection{Ordered factorization and covariant completion}
\label{chr-factorization}

A common difference-square formula would be invalid here: the
product of two equal-displacement differences is symmetric in
the bond components and would discard the odd-displacement,
skew-component part of \(q(r)\).  Even the forward-gradient Maxwell
kernel has such shifted entries.  Instead use an ordered Laurent
factorization.

For a Laurent series \(F\), define
\[
 {\cal B}_jF(z)=
 \frac{F(1,\ldots,1,z_j,\ldots,z_d)
       -F(1,\ldots,1,1,z_{j+1},\ldots,z_d)}{z_j-1}.
\]
The quotient is removable.  Applying the telescoping identity
twice, and using \eqref{chr-moments}, gives
\begin{equation}
 q(z)=\sum_{\rho,\tau}(z_\rho^{-1}-1)
       H_{\rho\tau}(z)(z_\tau-1),\qquad
 H_{\rho\tau}(z)=-z_\rho{\cal B}_\tau{\cal B}_\rho q(z).
 \label{chr-laurent}
\end{equation}
Each monomial division costs at most \(|r|_1\) terms with no
increase in exponent length.  Spending an exponential reserve
twice therefore yields a uniform exponentially summable \(H\).
For example an input exponent \(\sigma\) gives an output
exponent \(\sigma/2\) and a bound
\[
 \|H\|_{\sigma/2}
 \leq \frac{16e^{\sigma/2}}{e^2\sigma^2}\|q\|_\sigma
\]
in the combined absolute row norm; the analogous column bound
follows in the same way.  Replace \(H\) by \((H+H^T)/2\), where
the transpose includes displacement reversal and both gradient
and bond indices.  The factorization is unchanged.  We have
\begin{equation}
                  q=D_0^TH_0D_0,\qquad H_0=H_0^T.
 \label{chr-flat-factor}
\end{equation}
The color-scalar property is preserved.

Dress each \(H_0\) entry by its endpoint transport:
\[
 H^P(a)_{x\rho\alpha,y\tau\beta}
       =(H_0)_{\rho\alpha,\tau\beta}(y-x)W_{xy}(a).
\]
Let \(D_a\) be the nearest-neighbor covariant component gradient
and define
\begin{equation}
                  Q_{\mathrm{cov}}(a)
                   =D_a^TH^P(a)D_a .
 \label{chr-completion}
\end{equation}
This family is exactly covariant and transpose symmetric, with
\(Q_{\mathrm{cov}}(0)=q\).  Its first two Taylor kernels have
uniform exponentially rooted bounds: path derivatives only add
polynomial length factors.  On a sufficiently small complex
current-angle ball, the path-product factor
\(\exp(C\|a\|_\infty|r|)\) is paid by the exponential reserve.

Apply the local-gauge argument to \(Q_{\mathrm{cov}}\) as well
as to \(L\).  They have the same flat \(q\), hence the same dressed
flat comparison \(L^P\).  Subtraction of their two estimates
\eqref{chr-peierls-error} gives symmetric errors
\begin{equation}
 {\cal E}_j=[t^j](L(t\mu A)-Q_{\mathrm{cov}}(t\mu A)),
 \qquad
 \|{\cal E}_j\|_{\mathrm{Schur}}
           \leq C\mu^2a_*^j,\quad j=1,2.
 \label{chr-completion-error}
\end{equation}
This argument does not presume that the paths generated by
\eqref{chr-completion} coincide with the initially selected
endpoint paths.

\subsection{Coefficient expansion and localization}
\label{chr-expansion}

Write, through second order,
\[
 D_{t\mu A}=D_0+tF_1+t^2F_2,\qquad
 H^P(t\mu A)=H_0+tH_1+t^2H_2.
\]
Nearest-neighbor differentiation and exponential path moments give
\begin{align}
 \|F_1\|+\|H_1\|&\leq C\mu a_*,
 \notag\\
 \|F_2\|+\|H_2\|&\leq C\mu^2a_*^2 .
 \label{chr-factor-bounds}
\end{align}
All \(H_j\) are symmetric.  Set
\begin{align}
 Z_1&=\tfrac12H_1D_0+H_0F_1,
 \notag\\
 Z_2&=\tfrac12H_2D_0+H_1F_1+H_0F_2,
 \notag\\
 E_1&={\cal E}_1,\qquad
 E_2=F_1^TH_0F_1+{\cal E}_2 .
 \label{chr-exact-z}
\end{align}
Direct expansion of \eqref{chr-completion} proves
\eqref{chr-decomposition}.  The counting norm of \(D_0\) is bounded
by a dimensional constant, so \eqref{chr-factor-bounds} and
\eqref{chr-completion-error} prove \eqref{chr-bounds}.
Finally,
\[
 \|D_0R_0^{-1/2}\|\leq1,\qquad
 \|R_0^{-1/2}\|\leq\mu^{-1},
\]
which proves \eqref{chr-relative}.

For localization, use complete directions supported in a fixed
smooth physical partition with bounded overlap and bounded physical
first derivatives.  The terms \(F_1,F_2\) have nearest-neighbor
support.  Each marked source in \(H_1,H_2\) lies on a selected
endpoint path.  A distant source forces a long path, and the
exponential \(H_0\) kernel pays its length.  In the remainders,
a source on a local-gauge path is connected to the fixed matrix
row through the original background-source endpoint.  The
original rooted Taylor bound pays that distance.  Discrete
Stokes uses source differences in the same smooth collar,
enlarged by one fine edge.  Polynomial path multiplicities are
absorbed by a smaller exponential exponent.

Thus the factors in \eqref{chr-exact-z} and the \(E_j\) have
rooted final-block bounds of the form
\(\exp[-c\tau(\text{row},\text{column},\text{source labels})]\)
with the powers in \eqref{chr-bounds}.  Fine-coordinate Schur
summation proves the block-operator estimate; summing \(K^4\)
independent entry bounds is unnecessary.  No sharp source cutoff
is differentiated.

Physical weights \(\exp(\theta f(x/K))\), with \(f\) Lipschitz,
change a path factor by at most
\(\exp(C\theta\mu|r|)\), again paid by the exponential reserve.
Commuting \(D_0\) past the weight produces \(O(\theta\mu)\).
When multiplied by \(Z_1\) or \(Z_2\), this has the zero-order
size required for \(E_1\) or \(E_2\).  This supplies the weighted
localized form of the theorem.

\subsection{Scope of the actual application}
\label{chr-scope}

For \(L=S_M''\), the hypotheses are the uniform action support and
Cauchy bounds through degree four, actual gauge covariance, and
the already normalized flat kernel.  The result concerns bounded
physical \(C^1\) fields and their localized jets.  It does not give
the same gain for rapidly oscillating bounded potentials, whose
fine curvature may be \(O(K^{-1})\), and it gives no estimate for
the hard-projection scalar operator.

In the intended soft-vector application, the exact on-shell
identification and its source-coordinate conventions must first
justify the relevant affine jets of
\[
 S_M''(e^{iA/K})+D_AD_A^T+
 \overline C_A^T
 \bigl[aI-S_{MK}''(M_K(e^{iA/K}))\bigr]\overline C_A .
\]
The corrected constraint's full \(O(K^{-1})\) bound and its
physical source derivatives are separate inputs.  The present
theorem supplies the action-Hessian relative vertices.  It does
not by itself identify a soft-vector determinant, control every
stationary response, or determine a one-loop sign.



\section{A soft vector completion with a bounded coarse remainder}
\label{svc-section}

The exact identities below use the actual mean-zero determinant.
The large-$K$ estimate concerns its first two affine background jets.
Its auxiliary off-shell family retains the effective action $S_M$;
it is not replaced by a bare Wilson Hessian.

\subsection{Exact finite factorization}
\label{svc-factor-section}

Use current counting coordinates, $K=\ell^k$, $\mu=K^{-1}$, and
\begin{equation}
 J=K^2B^T,\qquad J^TJ=I,\qquad \Pi=JJ^T,\qquad P=I-\Pi.
 \label{svc-projectors}
\end{equation}
The actual covariant pair and residual Hessian are
\begin{equation}
 \widehat G=PD^TDP+\mu^2\Pi,\qquad
 Q_B=\mathcal H+DPD^T+a\bar C^T\bar C,
 \qquad \mathcal H=L-C^TqC.
 \label{svc-actual}
\end{equation}
Here $L=S_M''-\lambda\cdot f''$, $q=S_{MK}''$, all Hessians are
left-relative, and $a\ge R_q+1$ is fixed. The actual correction is
$\mathcal D=\mathcal N_0(\log\det\widehat G-\tfrac12\log\det Q_B)$.
Define its soft vector completion and coarse Schur factor by
\begin{equation}
 V=\mathcal H+DD^T+a\bar C^T\bar C
        =Q_B+(DJ)(DJ)^T,
 \qquad S=I-J^TD^TV^{-1}DJ.
 \label{svc-soft-vector}
\end{equation}
The determinant lemma and Woodbury identity give
\begin{equation}
 \det Q_B=\det V\det S,\qquad
 S=(I+J^TD^TQ_B^{-1}DJ)^{-1}.
 \label{svc-schur}
\end{equation}
These are identities between positive matrices on the real small germ,
with normalized holomorphic continuations through flatness.

For the soft scalar matrices
\begin{equation}
 T_s=D^TD+\mu^2\Pi,\qquad M_s=\mu^2J^TT_s^{-1}J,
 \label{svc-soft-scalar}
\end{equation}
the scalar Schur identity is $\det\widehat G=\det T_s\det M_s$.
Thus the exact actual decomposition is
\begin{equation}
 \boxed{\displaystyle
 \mathcal D=\mathcal N_0\left[
 \log\det T_s-\tfrac12\log\det V
                 +\log\det M_s-\tfrac12\log\det S\right].}
 \label{svc-four-factor}
\end{equation}
No massless based scalar factor or omitted Haar normalization occurs.
The matrix $S$ has one fixed color fiber per final scalar site and is
covariant in that coarse frame.

\subsection{Uniform flat spectra and energy resolvents}
\label{svc-flat-section}

A subscript $0$ denotes flat background. The proved mean-zero estimates
are $\Delta_1\le C_\Delta Q_{B,0}$ and
$Q_{B,0}\ge c\mu^2I$. Since
$D_0\Pi_0D_0^T\le D_0D_0^T\le\Delta_1$,
\begin{equation}
 Q_{B,0}\le V_0\le(1+C_\Delta)Q_{B,0},\qquad
 (1+C_\Delta)^{-1}I\le S_0\le I.
 \label{svc-flat-comparison}
\end{equation}
Indeed $\|Q_{B,0}^{-1/2}D_0J_0\|^2\le C_\Delta$ and
\eqref{svc-schur} give the second assertion. The number of coarse
rows has not been used as an estimate for this condition number.

At zero, $qD_c=0$, so the ordinary and corrected constraint terms
agree inside $C^TqC$. With $A_K=\bar C(0)$,
\begin{equation}
 V_0=q_M+D_0D_0^T+A_K^T(aI-q_{MK})A_K
       \ge\Delta_1+A_K^TA_K.
 \label{svc-flat-vector}
\end{equation}
The uniform analytic symbol of $q_M$ vanishes to order two at zero.
Its uniform second moments therefore give $q_M\le C\Delta_1$.
Also $\|A_K\|\le\mu$ and $\|q_{MK}\|\le R_q$.
Combining these facts with the gap gives
\begin{equation}
 c(\Delta_1+\mu^2I)\le V_0\le C(\Delta_1+\mu^2I),
 \label{svc-energy-comparison}
\end{equation}
on the chosen infinite structural system and its explicit periodizations. The finite actual lower bound, gap and heat estimate below remain valid without identifying arbitrary finite winding kernels with those periodizations.

For the spatial argument, $q_M$ is the kernel of the chosen infinite structural system. The same conclusions hold for its explicitly periodized kernels; they are not asserted for unrelated finite winding kernels merely by sampling. For later source estimates the spatial exponent must also be uniform
in final units. If $\delta(p)$ is the scalar lattice symbol, the common
analytic strip and double zero of $q_M$ imply
\begin{equation}
 \|q_M(p+i\mu\theta)-q_M(p)\|
 \le C\{\mu|\theta|\sqrt{\delta(p)}+\mu^2|\theta|^2\}.
 \label{svc-fine-twist}
\end{equation}
Near zero this is Taylor's formula; away from zero it follows from
bounded derivatives. The ordinary longitudinal term has the same bound.
Conjugating fine fields by $e^{\theta\cdot x/K}$ and coarse fields by
$e^{\theta\cdot y}$, straight-average row and column counting gives
\begin{equation}
 \|A_{K,\theta}\|\le C\mu,\qquad
 \|A_{K,\theta}-A_K\|\le C|\theta|\mu.
 \label{svc-average-twist}
\end{equation}
The coarse symbol $q_{MK}$ has a uniform strip. Its coarse correction
therefore changes by $O(|\theta|\mu^2)$.
Measured between two $(\Delta_1+\mu^2I)^{-1/2}$ factors, the total
change is $O(|\theta|)$. Equations \eqref{svc-energy-comparison}--
\eqref{svc-average-twist} give a common physical strip on which all
energy sandwiches of $V_0^{-1}$ are bounded uniformly. The energy legs
are the ordinary componentwise first differences and $\mu I$.

Bloch decomposition by $K$-blocks turns these into bounded
operator-valued analytic functions. Fourier integration gives uniform
exponentially summable final-block kernels. The internal fiber may
have dimension $K^4$, but the bounds are Hilbert operator norms; no
entrywise sum over that fiber is taken. Explicit periodizations obey the same bounds. The established fixed-$K$ local-limit interface identifies the resulting infinite source coefficient with the actual local Maxwell coefficient; it does not identify arbitrary finite kernels with these periodizations. In particular $D_0^TV_0^{-1}D_0$, $S_0$ and, after shrinking
the exponent using \eqref{svc-flat-comparison}, $S_0^{-1}$ have uniform
weighted block bounds.

Min--max also gives the flat heat estimate
\begin{equation}
 N^{-4}\operatorname{Tr}e^{-uK^2V_0}
 \le C\left\{\frac{K^4}{(1+K^2u)^2}+N^{-4}\right\}e^{-cu}.
 \label{svc-flat-heat}
\end{equation}
It follows from $V_0\ge Q_{B,0}$ and the mean-zero heat theorem.
The same Mellin integration gives bounded inverse moments for powers
$p>2$, the logarithmic bound at $p=2$, and exponential final-time tails.
These flat estimates alone say nothing about a magnetic derivative.

\subsection{The auxiliary family and the exact affine matching}
\label{svc-auxiliary-section}

For bounded physical $C^1$ potentials $A$ put $U_A=e^{i\mu A}$ and
\begin{align}
 \widetilde V(A)
  &=S_M''(U_A)+D_AD_A^T
       +\bar C_A^T\{aI-S_{MK}''(M_K(U_A))\}\bar C_A,
       \label{svc-auxiliary-vector}\\
 \widetilde S(A)
  &=I-J_A^TD_A^T\widetilde V(A)^{-1}D_AJ_A.
       \label{svc-auxiliary-coarse}
\end{align}
Every Hessian in these formulas is the actual relative-coordinate
classical Hessian. The family is covariant, but it is not declared equal
to \eqref{svc-soft-vector} off shell.

On the established commuting affine probe, the actual fine and coarse
currents and the multiplier vanish through the two jets used in the
Maxwell extraction. Thus $qD_c=0$ and
\begin{equation}
 C^TqC=\bar C^Tq\bar C,\qquad L=S_M''
 \quad\hbox{through the first two affine jets}.
 \label{svc-on-shell-match}
\end{equation}
The actual smooth stationary polynomial jets are the corresponding
fine affine connection. Hence $V,S$ and
$\widetilde V,\widetilde S$ have identical first two affine jets.
This is the precise interface used below.

\subsection{Relative vertices with localized sources}
\label{svc-vertices-section}

Use complete physical source fields with
\[
 a_* =\|A\|_\infty+
       \mu^{-1}\max_{\text{nearest, same component}}|A(e')-A(e)|.
\]
The corrected relative-Hessian theorem supplies, for $j=1,2$,
\begin{align}
 [t^j]S_M''(e^{it\mu A})
     &=d_0^TZ_j+Z_j^Td_0+E_j,\notag\\
 \|Z_1\|&\le C\mu a_*,&\quad \|E_1\|&\le C\mu^2a_*,\notag\\
 \|Z_2\|&\le C\mu^2a_*^2,&\quad \|E_2\|&\le C\mu^2a_*^2.
 \label{svc-relative-vertices}
\end{align}
Here $d_0$ is the componentwise gradient of a bond field. These are
source-rooted block estimates, including physical weights.
Their proof uses the ordered Laurent factorization
$q_M=d_0^TH_0d_0$ and a symmetric covariant completion. A common
componentwise-difference formula would be false for the shifted
component-skew kernel and is not used.

The first two $D_A$ jets are multiplication operators of sizes
$O(\mu a_*)$ and $O(\mu^2a_*^2)$. Their contribution to
$D_AD_A^T$ has the same form as \eqref{svc-relative-vertices}.
The actual moving-reference products obey
$\|\bar C_A\|=O(\mu)$ with first two physical-background derivatives
of the same scale. Their rows and sources lie in two final blocks.
The ordinary output $M_K(e^{i\mu A})$ is uniformly analytic and
$O(\|A\|)$ in final units, with the same fixed physical range.
The previously proved relative coarse Hessian kernels are uniform
in $MK$. Therefore the last term of
\eqref{svc-auxiliary-vector}, including its first two derivatives,
is a zero-order $O(\mu^2)$ vertex with exponentially rooted sources.
These conclusions use corrected moving tangents, not an estimate
for the full uncorrected ordinary adjoint.

The normalized block map $J_A$ and its first two physical derivatives
have uniform own-block operator bounds. Keep the factorization
\begin{equation}
 R(A)=D_A^T\widetilde V(A)^{-1}D_A,
 \qquad \widetilde S(A)=I-J_A^TR(A)J_A.
 \label{svc-avoid-interface}
\end{equation}
The derivatives of $J$ stay outside $R$. In particular, no raw derivative
of $DJ$ at a sharp block interface is estimated before using a resolvent.
Since $V_0\ge D_0D_0^T$ and $V_0\ge c\mu^2I$,
\begin{equation}
 \|V_0^{-1/2}D_0\|\le1,\qquad
 \|V_0^{-1/2}D'_0\|\le Ca_*.
 \label{svc-gauge-legs}
\end{equation}
For example,
\begin{equation}
 R'=D'^TV_0^{-1}D_0+D_0^TV_0^{-1}D'
           -D_0^TV_0^{-1}V'V_0^{-1}D_0.
 \label{svc-first-response}
\end{equation}
Differentiating once more yields the ordered two-$V'$ terms, one
$V''$ term, and the $D''$ and mixed $D',V'$ terms. Insert
\eqref{svc-relative-vertices} between the uniform weighted energy
resolvents. Each term is bounded with its physical source norms,
by \eqref{svc-gauge-legs}; this proves both first and second response
bounds without a fine-coordinate dimension factor.

For unit directions supported in smooth physical collars $x,y$, a
spare exponent attaches the sources to their resolvent legs and gives
\begin{align}
 \|\widetilde S'_x(i,j)\|&\le Ce^{-m\tau(i,j,x)},\notag\\
 \|\widetilde S''_{x,y}(i,j)\|
      &\le Ce^{-m\tau(i,j,x,y)}.
 \label{svc-rooted-sources}
\end{align}
Here $\tau$ is connecting-tree length including every displayed operator
endpoint and source. With a smaller exponent, these imply
$\sum_{i,j}e^{m'[d(i,x)+d(j,x)]}\|\widetilde S'_x(i,j)\|\le C$
and $\sum_{i,j}\|\widetilde S''_{x,y}(i,j)\|\le Ce^{-m'd(x,y)}$.
The full rooted estimate, rather than only these summed consequences,
is retained in the diagonal-density argument below.
The indices are final-block labels, with fixed color fibers.
The same source bounds hold after multiplying by $J,J',J''$, because
these are normalized own-block operators. Completeness of the fine
field collars is essential; sharp background cutoffs are not differentiated.

\subsection{The coarse logarithm and its actual Maxwell bound}
\label{svc-maxwell-section}

Equations \eqref{svc-flat-comparison} and
\eqref{svc-rooted-sources} give for the total Hessian
\begin{equation}
 K_S(x,y)=\operatorname{Tr}\left[
 S_0^{-1}\widetilde S''_{x,y}
 -S_0^{-1}\widetilde S'_xS_0^{-1}\widetilde S'_y\right],
 \qquad |K_S(x,y)|\le Ce^{-m' d(x,y)}.
 \label{svc-log-hessian}
\end{equation}
Both traces are coarse traces with fixed color dimension.

For the local diagonal density one must not cyclically rearrange a
partial trace. Put $R_t=(S_0+tI)^{-1}$. Its second source derivative
at root zero is the partial trace of the integral over $t\ge0$ of
\begin{equation}
 R_t\widetilde S''_{x,y}R_t
 -R_t\widetilde S'_xR_t\widetilde S'_yR_t
 -R_t\widetilde S'_yR_t\widetilde S'_xR_t.
 \label{svc-density-resolvent}
\end{equation}
The weighted bound is $\|R_t\|_m\le C/(1+t)$.
Thus its rooted density kernel is bounded by
$Ce^{-m''\tau(0,x,y)}$. This gives uniform polynomial source moments,
not merely decay in the source separation.

Covariance makes each diagonal color trace invariant. Apply the
rooted bounds to a smooth physical partition of an affine potential;
the collar norms grow only linearly with their distance from zero and
are summable against the rooted exponential. The density Ward and
reanchoring identity then yields
\begin{equation}
 \left|\beta\left(\tfrac12\mathcal N_0\log\det S\right)\right|
       \le C_S,
 \label{svc-bounded-remainder}
\end{equation}
uniformly in $M,K$, in the actual joint infinite structural systems
and polynomial-jet convention. The transfer from $\widetilde S$ to the
actual $S$ uses exactly \eqref{svc-on-shell-match}. The linear color jet
vanishes by semisimple conjugation invariance.

Combining \eqref{svc-four-factor} with the scalar coarse theorem gives
\begin{equation}
 \beta(\mathcal D)
  =\beta\!\left(\mathcal N_0
       [\log\det T_s-\tfrac12\log\det V]\right)+O(1),
 \label{svc-final-reduction}
\end{equation}
with a remainder uniform in old depth and endpoint. For the vector
coefficient, the auxiliary expression
\eqref{svc-auxiliary-vector} is valid through its two affine jets.
The fine soft pair has not yet been compared to a standard massive
vector/ghost reference. A fixed-integrand universality theorem alone
does not supply bounds uniform in this family. This module proves
neither that remaining logarithmic comparison nor its sign, and it
does not assert a full $K$-independent nonlinear background domain.
Finite-volume winding and normalized-extraction issues retain their
separate regulator hypotheses.



\section{The fine short-time band of the soft pair}
\label{suv-section}

This argument is on the chosen infinite structural system, in current
counting coordinates, with $\mu=K^{-1}$. All unbounded backgrounds are
formal polynomial jets of the commuting affine probe
\[
 U_\alpha(t;x)=\exp(i t x_\nu H/K^2),\qquad \alpha\ne\nu.
\]
It asserts no continuously varying magnetic flux on a finite torus.
The actual action $S_M$ is retained.

Put
\begin{align}
 B_s(t)&=D(t)^TD(t),&
 B_v(t)&=S_M''(U(t))+D(t)D(t)^T,\notag\\
 T_s(t)&=B_s(t)+\mu^2J(t)J(t)^T,&
 \widetilde V(t)&=B_v(t)+\bar C(t)^TF(t)\bar C(t),\notag\\
 F(t)&=aI-S_{MK}''(M_K(U(t))).
 \label{suv-operators}
\end{align}
The scalar and vector final-block-zero projectors are denoted
$P_{0,s}$ and $P_{0,v}$.

\begin{theorem}
\label{suv-main}
There is a constant independent of old depth $M$ and endpoint $K$ such
that
\begin{align}
 \int_0^1\left|[t^2]\operatorname{Tr}
       (P_{0,s}e^{-sT_s(t)})\right|\frac{ds}{s}&\le C,\notag\\
 \int_0^1\left|[t^2]\operatorname{Tr}
       (P_{0,v}e^{-s\widetilde V(t)})\right|\frac{ds}{s}&\le C.
 \label{suv-bounds}
\end{align}
The vector expression equals the actual soft vector through its two
affine jets. The effect of each soft coarse penalty alone is bounded
by $C\mu^2s$ before integration.
\end{theorem}

\subsection{Rooted estimates for the base operators}
\label{suv-base-section}

At flatness $B_s$ and $B_v$ have uniform exponential fine-current
Schur bounds. The vector statement uses the actual $q_M$; its first
two background derivatives use the uniform action kernels through
degree four and exact left-relative covariance. The gauge terms
are nearest-neighbor. The semigroup power series therefore has
Schur bound $Ce^{Cs}$ for $0\le s\le1$.

At a diagonal fine root $x$, a Cartan gauge subtracts the constant
potential at that root. The resulting source is
$t(z_\nu-x_\nu)H/K^2$. Diagonal color trace is unchanged.
Each first magnetic insertion now supplies $K^{-2}$ times a polynomial
in distances to the root, and each second insertion supplies $K^{-4}$.
The rooted exponential kernels pay the polynomial factors. The ordered
one-insertion and two-insertion heat formulas give, for $E=B_s,B_v$,
\begin{equation}
 \left|[t^2]\operatorname{tr}_{\mathfrak g}
      (e^{-sE(t)})_{xx}\right|
       \le CK^{-4}(s+s^2)e^{Cs},\qquad 0<s\le1.
 \label{suv-fine-root}
\end{equation}
There are $K^4$ scalar roots or $dK^4$ initial bond roots in one block.
Thus
\begin{equation}
 \left|[t^2]\operatorname{Tr}(P_{0,E}e^{-sE(t)})\right|
       \le C(s+s^2)e^{Cs}.
 \label{suv-base-bound}
\end{equation}
This proof uses no spectral gap. Polynomial gauges are justified by
compact truncation and exponential absolute moments, not by evaluating
an unbounded field inside a global analytic ball.

\subsection{Fixed coarse fibers for the penalty insertion}
\label{suv-penalty-section}

Write either perturbation as $C(t)^TF(t)C(t)$, with
\begin{equation}
 (C,F)=
 \begin{cases}
   (\mu J^T,I),&\text{scalar},\\
   (\bar C,aI-S_{MK}''(M_K(U))),&\text{vector}.
 \end{cases}
 \label{suv-factor-table}
\end{equation}
The output fiber of $C$ has fixed color/site or color/bond dimension
per final block. The actual moving-reference physical estimates,
centered on each dependency collar, give for $j=0,1,2$
\begin{equation}
 \|C_j(Y,X)\|\le C\mu(1+|X|+|Y|)^p
                         \mathbf1_{d(X,Y)\le r_0}.
 \label{suv-C-jets}
\end{equation}
Here $C_j=[t^j]C(t)$, and $r_0,p$ are fixed. For $J$ the support
is its own block. Centering the affine field uses bounded physical
$C^1$ directions; restoring the gauge frames costs only a fixed
polynomial in the block labels.

The ordinary map sends the affine probe to the coarse affine
connection. The actual coarse Hessian source bounds consequently give
\begin{equation}
 \|F_j(Y,Z)\|\le C(1+|Y|+|Z|)^p e^{-\kappa d(Y,Z)}.
 \label{suv-F-jets}
\end{equation}
These are whole-block Hilbert operator bounds, with no sum over
$K^4$ independent fine coordinates.

The base jets also have polynomially weighted exponential final-block
operator kernels. Neighboring blocks are controlled by the fine Schur
bound; separated blocks gain the fine exponential at distance of order
$K$ times the block distance. The flat semigroup series and its first
two insertion formulas therefore give the same kind of bound,
uniformly for $0\le s\le1$, for both $B$ and
$W=B+C^TFC$. This step takes no fine trace.

The exact Duhamel formula is
\begin{equation}
 e^{-sW(t)}-e^{-sB(t)}
   =-\int_0^s e^{-(s-u)W(t)}C(t)^TF(t)C(t)e^{-uB(t)}\,du.
 \label{suv-duhamel}
\end{equation}
Keep the fine root projector in the trace:
\begin{equation}
 \operatorname{Tr}_f[P_0 E_W C^TFC E_B]
       =\operatorname{Tr}_c[FC E_B P_0 E_W C^T].
 \label{suv-cyclic}
\end{equation}
The right side sums all coarse diagonal blocks; it is not a single
coarse block-zero density.

Expand its degree-two coefficient in the finitely many ordered ways.
There are always two $C$ factors, supplying $\mu^2$.
Equations \eqref{suv-C-jets}--\eqref{suv-F-jets} and the semigroup
bounds connect both coarse endpoints to block zero. Their absolute
sum is bounded by
\[
 C\mu^2\sum_{Y,Z}(1+|Y|+|Z|)^p
       e^{-\kappa'(|Y|+|Z|+d(Y,Z))}\le C'\mu^2.
\]
The trace is taken in a fixed-dimensional coarse fiber. A product
through the large fine block Hilbert space costs this coarse output
rank, not the fine dimension. Integrating \eqref{suv-duhamel} gives
\begin{equation}
 \left|[t^2]\operatorname{Tr}
       P_0(e^{-sW(t)}-e^{-sB(t)})\right|
          \le C\mu^2s,\qquad 0<s\le1.
 \label{suv-correction-bound}
\end{equation}
For compact source truncations, \eqref{suv-cyclic} is ordinary
finite-rank trace cyclicity. The displayed rooted sum supplies its
infinite limit and permits coefficient extraction. This also justifies
the complex-bilinear continuation at the real flat expansion point.

Combining \eqref{suv-base-bound} and \eqref{suv-correction-bound},
dividing by $s$, and integrating proves \eqref{suv-bounds}.
The exact affine matching of the soft vector transfers this estimate
to the actual vector coefficient.

\subsection{Scope}
\label{suv-scope}

The band $0<s<1$ is $0<u<K^{-2}$ in final physical heat time.
Its contribution to the actual signed soft-pair coefficient is bounded
uniformly. This proof does not compute the logarithmic coefficient on
$1<s<K^2$, compare the actual effective action to a standard massive
vector/ghost reference, or establish a large-time cancellation.
Finite winding and regulator identification retain their separate
hypotheses.


\section{A uniform affine infrared tail from local cubic resolvents}
\label{sft-section}

This section bounds the large-physical-time contribution of the actual
soft scalar/vector pair. The proof retains the root block in every
trace and introduces no factor counting its \(K^4\) fine variables.
It uses the localized relative vertices proved above, not only a
global relative-form estimate. It does not determine the
intermediate-time logarithmic coefficient.

\subsection{Precise operator and source inputs}
\label{sft-input-section}

Write \(\mu=K^{-1}\). Fine scalar or positive-bond spaces use counting
inner products, with fixed color/component multiplicities. Assign a
bond to the block containing its initial site. Let \(P_x\) be the
orthogonal projection onto the final block \(x\), and put
\begin{equation}
 \mathcal L=I+K^2\Delta,
 \label{sft-reference}
\end{equation}
where \(\Delta\) is the ordinary componentwise flat lattice Laplacian.
Let \(X_0\) denote the flat physical operator, namely \(K^2\) times
the corresponding counting matrix.

For the vector spatial estimates, work on the chosen joint infinite
structural system or on explicit periodizations of its kernels.
Actual finite-regulator gaps and flat heat estimates are separate
facts. A finite momentum grid does not imply exact infinite moment
cancellation, and arbitrary finite winding kernels are not identified
with these periodizations. The following abstract proof is also
finite-volume uniform whenever its three inputs hold on the finite
system in question.

The first input is uniform positive ellipticity:
\begin{equation}
 X_0=X_0^*,\qquad X_0\ge c_0\mathcal L.
 \label{sft-ellipticity}
\end{equation}
The second is a relative physical-weight estimate. For every bounded
physical \(1\)-Lipschitz function \(f\), set \(W=e^{\theta f}\),
using the fine-site position divided by \(K\). On bonds use their
initial sites. For sufficiently small fixed \(\theta>0\), assume
\begin{equation}
 \left\|\mathcal L^{-1/2}
       (WX_0W^{-1}-X_0)\mathcal L^{-1/2}\right\|
       \le C\theta .
 \label{sft-weighted-input}
\end{equation}
Bounded truncations of distance functions suffice.

For the third input, let \(a_j\) be a complete physical \(C^1\) field
supported in a fixed smooth collar of block \(j\), and write
\begin{equation}
 a_j^*=\|a_j\|_{\mathrm{phys},1}
 =\|a_j\|_\infty+
 K\max_{\substack{\mathrm{nearest}\\\mathrm{same\ component}}}
       |a_j(e')-a_j(e)|.
 \label{sft-source-norm}
\end{equation}
Denote the first and mixed second variations of \(X\) at zero by
\(V_j\) and \(V_{jk}\). For every weight allowed above, require
\begin{align}
 \|\mathcal L^{-1/2}WV_jW^{-1}\mathcal L^{-1/2}\|
     &\le Ca_j^*,\notag\\
 \|\mathcal L^{-1/2}WV_{jk}W^{-1}\mathcal L^{-1/2}\|
     &\le Ca_j^*a_k^*.
 \label{sft-conjugated-vertices}
\end{align}
In addition, let \(f_j\) be physical distance to the fixed physical
cube or collar of block \(j\), and \(W_j=e^{\theta f_j}\), initially
with bounded truncations. The marked-source bounds are
\begin{align}
 \|\mathcal L^{-1/2}W_jV_jW_j\mathcal L^{-1/2}\|
     &\le Ca_j^*,\notag\\
 \|\mathcal L^{-1/2}W_jV_{jk}W_j\mathcal L^{-1/2}\|
     &\le Ca_j^*a_k^*,
 \label{sft-marked-vertices}
\end{align}
with the second bound also centered at \(k\).
The function \(f_j\) is Lipschitz on fine physical positions:
it is not a piecewise constant block-index distance. Its infimum
and supremum on any one block differ by a dimensional constant.

The two-sided weights in \eqref{sft-marked-vertices} attach a source
to both operator ends. An unlocalized bound on
\(X_0^{-1/2}V_jX_0^{-1/2}\) would not supply this input.
All source jets may be holomorphic and bilinear. Only the flat
operators \(X_0,\mathcal L\) are required to be positive.

\subsection{Why the actual soft pair supplies these inputs}
\label{sft-application-inputs}

For \(T_s=D_A^TD_A+\mu^2\Pi_A\), the soft scalar theorem gives
\eqref{sft-ellipticity}. Its weighted \(D,J\) expansion gives
\eqref{sft-weighted-input}. Its first two source derivatives give
\eqref{sft-conjugated-vertices}--\eqref{sft-marked-vertices}.
Here \(J\) is normalized and own-block local, so its derivatives
have no fine-fiber dimension loss. Scalar sources even need only
their supremum norm.

For the soft vector, the mean-zero Laplacian domination and gap,
preserved by the positive \(DJ\) update, imply
\eqref{sft-ellipticity} for \(X_0=K^2V_0\).
The relative-Hessian theorem gives, after physical rescaling,
first-two source jets of the form
\begin{equation}
 D_{\mathrm{phys}}^TZ+Z^TD_{\mathrm{phys}}+E,
 \label{sft-energy-vertex}
\end{equation}
with uniformly bounded, exponentially source-rooted factors \(Z,E\).
The second-order \(Z\) may have an additional beneficial factor
\(\mu\). Since
\(\|D_{\mathrm{phys}}\mathcal L^{-1/2}\|\le1\),
these factors give the relative vertex bounds. For example,
\[
 W_jD_{\mathrm{phys}}^TZW_j
 =D_{\mathrm{phys}}^T(W_jZW_j)
   +[W_j,D_{\mathrm{phys}}^T]W_j^{-1}(W_jZW_j).
\]
The commutator in the second term is \(O(\theta)\); source-rooted
decay bounds \(W_jZW_j\) at a smaller exponent. This proves the
two-sided normalization without commuting a fractional power of
\(\mathcal L\) through a weight.

The covariant \(D_AD_A^T\) term has the same elementary structure.
Each corrected angle average \(\bar C_A\), and each of its first
two physical-background derivatives, has counting size \(O(\mu)\).
Their pair therefore gives \(O(\mu^2)\), canceled by \(K^2\).
The coarse Hessian has fixed color fiber and uniform source bounds.
These terms supply zero-order vertices in
\eqref{sft-energy-vertex}. The soft-vector module separately
identifies this auxiliary family with the actual stationary affine
jets through order two, retaining the current and residual terms.

For completeness, the action contribution to
\eqref{sft-weighted-input} uses precisely the repaired ordered
Laurent factorization
\begin{equation}
 q_M=d_0^TH_0d_0.
 \label{sft-ordered-factorization}
\end{equation}
Both the constant and linear flat momentum jets vanish, and \(H_0\)
is symmetric with uniformly exponentially summable current kernels.
Separate component row sums alone would not justify this step;
no common componentwise difference-square identity is used.
Conjugation changes each outer gradient by a bounded multiplication
operator \(O(\theta\mu)\), and changes \(H_0\) by \(O(\theta\mu)\)
using its exponential first moment. Keep the latter change between
the two gradients. After multiplication by \(K^2\) and the two
\(\mathcal L^{-1/2}\) factors, the outer errors are
\(O(\theta+\theta^2)\) and the middle error is \(O(\theta\mu)\).
The longitudinal term obeys the same estimate. The bounded physical
block correction has fixed physical range, with a coarse Hessian
having uniform exponential moments, so its error is \(O(\theta)\).
This proves \eqref{sft-weighted-input} without a false uniform bound
on the absolute row norm of \(X_0\), which can grow as \(K^2\).

The uniform claims concern the explicit auxiliary source estimates.
Their identification with the actual affine coefficient is used at
each fixed \(K\); no uniform stationary-response bound is needed.

\subsection{A block Schatten bound without a rank loss}
\label{sft-schatten-section}

The elementary massive lattice heat estimate gives
\begin{equation}
 \operatorname{Tr}(P_x\mathcal L^{-3}P_x)\le C
 \label{sft-third-moment}
\end{equation}
uniformly in \(K\) and compatible volume. Indeed, integrate the
flat heat trace against \(t^2e^{-t}/2\). On a torus of current side
\(KN\), the heat trace per final block is at most
\[
 C\left\{\frac{K^4}{(1+K^2t)^2}+N^{-4}\right\}.
\]
Its indicated integral is bounded, including the finite zero mode.
The infinite lattice obeys the same bound without that mode.

Compression is not interchanged with taking powers. If \(u_i\) is
an orthonormal eigenbasis of \(P_x\mathcal L^{-1}P_x\) in the block,
scalar Jensen applied to the spectral measure of \(\mathcal L^{-1}\)
gives
\(\langle u_i,\mathcal L^{-1}u_i\rangle^3
 \le\langle u_i,\mathcal L^{-3}u_i\rangle\).
Therefore
\begin{equation}
 \|P_x\mathcal L^{-1/2}\|_{\mathfrak S_6}^6
 =\operatorname{Tr}[(P_x\mathcal L^{-1}P_x)^3]
 \le\operatorname{Tr}(P_x\mathcal L^{-3}P_x)\le C.
 \label{sft-s6}
\end{equation}
The analogous right-ended norm has the same bound. Schatten
Hölder with exponents \(6,\infty,6\) consequently yields
\begin{equation}
 \|P_x\mathcal L^{-1/2}A\mathcal L^{-1/2}P_y\|_{\mathfrak S_3}
 \le C\|A\|.
 \label{sft-s3}
\end{equation}
There is no factor proportional to the block rank \(K^4\).

\subsection{A common contour and rooted resolvent derivatives}
\label{sft-resolvent-section}

Choose \(0<\lambda_0<c_0/4\) and \(0<\varphi<\pi/2\).
Let \(\Gamma\) be the positively oriented Dunford contour on the
rays \(\lambda_0+r e^{\pm i\varphi}\), enclosing the positive
spectrum to their right. At finite regulator, first close it beyond
the largest eigenvalue. For \(R_z=(z-X_0)^{-1}\),
\begin{equation}
 \|\mathcal L^{1/2}R_z\mathcal L^{1/2}\|\le C_\Gamma,
 \qquad z\in\Gamma .
 \label{sft-contour-bound}
\end{equation}
Indeed, \(X_0\ge c_0\mathcal L\), while the spectral norm of
\(X_0^{1/2}R_zX_0^{1/2}\) is bounded on those rays.
By \eqref{sft-weighted-input} and a relative resolvent Neumann
expansion, choosing \(\theta\) sufficiently small gives
\begin{equation}
 WR_zW^{-1}=\mathcal L^{-1/2}A_{z,W}\mathcal L^{-1/2},
 \qquad \|A_{z,W}\|\le C ,
 \label{sft-twisted-resolvent}
\end{equation}
uniformly on the whole contour.

The first and mixed second derivatives are
\begin{align}
 \delta_jR_z&=R_zV_jR_z,\notag\\
 \delta_j\delta_kR_z
  &=R_zV_{jk}R_z+R_zV_jR_zV_kR_z+R_zV_kR_zV_jR_z.
 \label{sft-resolvent-derivatives}
\end{align}
Relative normalization cancels adjacent powers of \(\mathcal L\):
each product is \(\mathcal L^{-1/2}A\mathcal L^{-1/2}\)
with uniformly bounded \(A\).

Source localization can be checked without any fractional-weight
commutator. For a marked source \(j\),
\begin{equation}
 W_j(R_zV_jR_z)W_j
 =(W_jR_zW_j^{-1})(W_jV_jW_j)(W_j^{-1}R_zW_j).
 \label{sft-marked-first}
\end{equation}
For the ordered term with another source \(k\), put
\(R_+=W_jR_zW_j^{-1}\), \(R_-=W_j^{-1}R_zW_j\); then
\begin{equation}
 W_j(R_zV_jR_zV_kR_z)W_j
 =R_+(W_jV_jW_j)R_-(W_j^{-1}V_kW_j)R_-.
 \label{sft-marked-second}
\end{equation}
The marked vertex uses \eqref{sft-marked-vertices} and the other
vertex uses \eqref{sft-conjugated-vertices}. The primitive
\(V_{jk}\) has the marked bound at either source. Thus
\eqref{sft-s3} applies to every term.

Write \(\tau\) for the shortest connected graph length in the
final-block graph, allowing auxiliary vertices. The resulting
block bounds, for \(n=0,1,2\), are
\begin{equation}
 \|P_x\delta_{j_1}\cdots\delta_{j_n}R_zP_y\|_{\mathfrak S_3}
 \le C_n\prod_{\ell=1}^n a_{j_\ell}^*
       e^{-c\tau(x,y,j_1,\ldots,j_n)}.
 \label{sft-rooted-resolvent}
\end{equation}
For \(n=0\), use a one-sided distance weight from \(y\).
For \(n=1\), the marked weight gives
\(e^{-\theta[d(x,j)+d(y,j)]}\). For \(n=2\), take the
geometric mean of the separate \(j\)- and \(k\)-centered estimates.
Their four edges connect all marked vertices and dominate a
connecting graph length. The infimum/supremum comparison of
physical distances within each block changes only the constant.
Fixed collars have the same effect. Periodic physical distances
and bounded truncations give the identical conclusion whenever the
inputs hold; the uniform bounds permit removal of the truncations.

\subsection{Three resolvents preserve the local trace}
\label{sft-heat-section}

Functional calculus gives, for \(u>0\),
\begin{equation}
 e^{-uX}=\frac{2}{u^2}\frac1{2\pi i}
       \int_\Gamma e^{-uz}(z-X)^{-3}\,dz .
 \label{sft-cubic-calculus}
\end{equation}
The scalar residue of the third-order pole is
\(u^2e^{-u\lambda}/2\), fixing both coefficient and sign.

Differentiate \(R_z^3\) through order two. Each term is a product of
three differentiated resolvents, with derivative orders
\(n_1+n_2+n_3\le2\). Inserting the block resolution of the identity
between them gives
\begin{align}
 &\left|\operatorname{Tr}
  (P_0R_z^{(n_1)}P_iR_z^{(n_2)}P_jR_z^{(n_3)}P_0)\right|
 \notag\\
 &\qquad\le
  \|P_0R_z^{(n_1)}P_i\|_{\mathfrak S_3}
  \|P_iR_z^{(n_2)}P_j\|_{\mathfrak S_3}
  \|P_jR_z^{(n_3)}P_0\|_{\mathfrak S_3}.
 \label{sft-local-holder}
\end{align}
Here superscripts denote the indicated source derivatives, not
ordinary powers. The root block stays in place. Nothing is moved
cyclically through \(P_0\).

Equation \eqref{sft-rooted-resolvent} and exponential convolution,
with a smaller exponent, imply
\begin{equation}
 \left|\delta_{j_1}\cdots\delta_{j_n}
       \operatorname{Tr}(P_0R_z^3P_0)\right|
 \le C_n\prod_{\ell=1}^n a_{j_\ell}^*
        e^{-c'\tau(0,j_1,\ldots,j_n)},\qquad n\le2.
 \label{sft-cubic-root}
\end{equation}
The two intermediate block sums are absolutely convergent in trace
norm, uniformly in \(K,z\) and the applicable volume. This constructs
the infinite local trace without an infinite extensive trace.

On the contour,
\(\operatorname{Re}z=\lambda_0+r\cos\varphi\), whence
\[
 \int_\Gamma |e^{-uz}|\,|dz|
       \le C u^{-1}e^{-\lambda_0u}.
\]
Combining this with
\eqref{sft-cubic-calculus}--\eqref{sft-cubic-root} proves
\begin{equation}
 \left|\delta_{j_1}\cdots\delta_{j_n}
       \operatorname{Tr}(P_0e^{-uX}P_0)\right|
 \le C_nu^{-3}e^{-\lambda_0u}
       \prod_{\ell=1}^n a_{j_\ell}^*
       e^{-c'\tau(0,j_1,\ldots,j_n)},\quad u\ge1,\ n\le2.
 \label{sft-rooted-heat}
\end{equation}
The displayed absolute majorants justify the differentiations,
contour integral, block sums and local limits. At fixed \(K\), begin
with bounded compact sources in the actual local germ, then extend
their coefficients using these bounds. No small-time asymptotic is
claimed by the harmless factor \(u^{-3}\).

\subsection{The actual affine infrared contribution}
\label{sft-affine-section}

Choose a smooth physical partition of unity on fixed block collars,
with uniformly bounded first derivatives and finite overlap.
Decompose the affine physical field
\begin{equation}
 A_\alpha(x)=\frac{x_\beta}{K}H,\qquad \alpha\ne\beta,
 \label{sft-affine}
\end{equation}
into complete supported fields \(a_j\).
Their norms satisfy \(a_j^*\le C(1+|j|)\). The second source bound
\eqref{sft-rooted-heat} is summable with these polynomial weights.
Consequently
\begin{equation}
 \left|[t^2]\operatorname{Tr}_0 e^{-uX(tA)}\right|
       \le Ce^{-cu},\qquad u\ge1,
 \label{sft-affine-heat}
\end{equation}
uniformly in \(K\) and old depth on the stated structural systems.
This is evaluation of exponentially localized Taylor kernels,
not insertion of a nonzero unbounded field into a fixed analytic ball.

The actual first two stationary affine jets agree with the explicit
representatives just estimated, up to a fine gauge. Their diagonal
color traces are individually gauge invariant. Thus the fixed-\(K\)
polynomial-jet identification transfers
\eqref{sft-affine-heat} to the actual coefficient before taking the
uniform bound. It does not require a uniform-\(K\) norm for the
full stationary response.

For \(X_s=K^2T_s\) and \(X_v=K^2V\), we obtain
\begin{theorem}[Uniform actual infrared tail]
\label{sft-main}
Under the established soft-factor and localized relative-vertex
inputs, the actual affine coefficient satisfies
\begin{equation}
 \int_1^\infty
  \left|[t^2]\left\{
   \operatorname{Tr}_{0,E}e^{-uX_v(t)}
     -2\operatorname{Tr}_{0,S}e^{-uX_s(t)}
                 \right\}\right|\frac{du}{u}\le C .
 \label{sft-tail}
\end{equation}
The constant is uniform in \(K\), old depth and the chosen joint
infinite structural systems. Each factor separately obeys the same
type of bound.
\end{theorem}

The fine link angle corresponding to \eqref{sft-affine} is
\(t x_\beta H/K^2\). The change of variables \(s=K^2u\)
identifies \eqref{sft-tail} with the current-time tail \(s\ge K^2\)
of the actual soft-pair Maxwell heat formula.
Constant determinant normalizations have zero affine coefficient.

No finite torus affine-flux family is asserted. Finite source bounds
apply to explicit periodizations or to finite models for which the
three inputs were separately verified. The actual polynomial
coefficient is obtained in the chosen infinite structural system by
the fixed-\(K\) identification; unrelated actual finite winding
kernels are not silently substituted.

Together with the exact soft factorization and the bounded coarse
Maxwell remainders, \eqref{sft-tail} removes the large-physical-time
part from the logarithmic calculation. The remaining comparison is
the short-current-time contribution and the intermediate range
\(1\ll s\ll K^2\). Neither the value nor the sign of the logarithmic
coefficient follows from this tail estimate.


\section{The next comparison}

The exact soft factorization and bounded coarse corrections reduce
the required logarithm to the fine soft scalar/vector pair. Its
short-current-time interval \(0<s<1\) and its physical infrared
tail \(s>K^2\) both have uniformly bounded affine contributions.
Writing
\[
 h^{\rm soft}_{M,K}(s)=[t^2]\left\{
   \operatorname{Tr}_{0,E}e^{-sV(t)}
       -2\operatorname{Tr}_{0,S}e^{-sT_s(t)}\right\},
\]
we therefore have the actual reduction
\begin{equation}
 \boxed{\displaystyle
 \beta(\mathcal D_{M,K})
       =\int_1^{K^2}h^{\rm soft}_{M,K}(s)\,\frac{ds}{s}+O(1),}
 \label{soft-final-band-reduction}
\end{equation}
uniformly in preceding depth and endpoint on the stated structural
systems. For \(K=1\) the intermediate interval is empty. The
growing hard-projection vertices have been removed from this
calculation by exact algebra. The actual effective-action Hessian
remains present, with its now controlled relative physical vertices.

Two distinct ingredients remain in identifying the logarithm.
First, the soft infrared penalties must be compared to a specified
massive reference without discarding their gauge cancellations.
Second, the actual effective-action vertices must be compared to
the reference vertices with constants uniform in preceding depth.
The previously proved nonlinear Ward extraction matches their
marginal momentum jets through field degree four; a summable
remainder estimate for the determinant is still required.
A sufficient concrete target, with \(\delta>0\), is
\[
 \left|h^{\rm soft}_{M,K}(s)-11\alpha_g\right|
       \le C\left(s^{-\delta}+(s/K^2)^\delta\right),
       \qquad 1\le s\le K^2.
\]
Its logarithmic-time integral gives the required
\(22\alpha_g\log K+O(1)\). This displayed comparison is a target,
not an estimate proved in this notebook.

For a fixed logarithmically divergent one-loop lattice integral,
Adams--Lee prove a universal logarithmic term with a finite remainder.
Their theorem provides a useful reference-integral endpoint once
its hypotheses and normalization are matched. It does not by itself
apply uniformly to our growing block fibers and varying effective
actions. The model-specific comparison must establish that interface.

A positive actual block drift would activate the finite transfer
and tuning results in Notebook 32. The full existence argument also
requires the recorded measure matching, mixed and large-field
histories, sourced continuum limits, nontriviality, Euclidean
covariance, positivity and reconstruction. These remain part of the
original goal. The spectral gaps here are Gaussian fluctuation
estimates, not a Yang--Mills mass-gap theorem.

\begin{thebibliography}{9}
\bibitem{priorsoft}
\emph{Toward a four-dimensional Yang--Mills construction},
Notebooks 32--35 and the earlier classical support and Ward
extraction modules in this workspace.
\bibitem{adamslee}
D.~H.~Adams and W.~Lee,
\emph{Structure of logarithmically divergent one-loop lattice
Feynman integrals}, Phys.\ Rev.\ D \textbf{77} (2008), 045010,
\href{https://arxiv.org/abs/0709.0781}{arXiv:0709.0781}.
\bibitem{sourcesoft}
E.~Sussman and collaborators,
\href{https://github.com/ethanwsussman/Yang--Mills-LLM-experiment}
{Yang--Mills-LLM-experiment}, deposited accounts pinned at
commit \texttt{e392ef4501\allowbreak 3a581664ba\allowbreak b76b908671\allowbreak 70b30fe10e}.
\end{thebibliography}
\end{document}
